\documentclass[11pt,a4paper]{article}
\usepackage[T1]{fontenc}
\usepackage[utf8]{inputenc}
\usepackage[margin=1in]{geometry}
\usepackage{amsmath,amssymb,amsthm}
\usepackage{stmaryrd}
\usepackage{booktabs}
\usepackage{tabularx}
\usepackage{array}
\usepackage{enumitem}
\usepackage{listings}
\usepackage{xcolor}
\usepackage[expansion=false]{microtype}
\usepackage{tikz}
\usetikzlibrary{arrows.meta,positioning,fit,calc}
\usepackage[htt]{hyphenat}
\usepackage[hidelinks,pdftitle={F1R3Comb-Mat v0.5}]{hyperref}
\setlength{\emergencystretch}{3em}

\newcommand{\code}[1]{\texttt{#1}}
\newcommand{\Comb}{\textsc{F1R3Comb}}
\newcommand{\Mat}{\textsc{F1R3Comb-Mat}}
\newcommand{\Kind}{\textsc{K1ndl1ng}}
\newcommand{\Camp}{\textsc{CampF1R3}}
\newcommand{\Bell}{\textsc{B3ll0ws}}
\newcommand{\Fx}{\textsc{F1R3x}}
\newcommand{\Rhobots}{\textsc{Rhobots}}
\newcommand{\MUST}{\textsc{must}}
\newcommand{\MUSTNOT}{\textsc{must not}}
\newcommand{\SHOULD}{\textsc{should}}
\newcommand{\SHOULDNOT}{\textsc{should not}}
\newcommand{\MAY}{\textsc{may}}

\newcommand{\quo}[1]{@#1}
\newcommand{\drp}[1]{{*}#1}
\newcommand{\nil}{\mathbf{0}}
\newcommand{\nequiv}{\equiv_{\mathcal{N}}}
\newcommand{\nms}[1]{\mathcal{N}^{*}(#1)}
\newcommand{\Nsort}{\mathsf{N}}
\newcommand{\Psort}{\mathsf{P}}
\newcommand{\red}{\longrightarrow}
\newcommand{\wred}{\longrightarrow^{*}}
\newcommand{\scong}{\equiv}
\newcommand{\comps}[1]{\lvert #1 \rvert}

\newcommand{\mm}{\mathsf{m}}
\newcommand{\dd}{\mathsf{d}}
\newcommand{\kk}{\mathsf{k}}
\newcommand{\fw}{\mathsf{fw}}
\newcommand{\bl}{\mathsf{b}_{\mathsf{l}}}
\newcommand{\br}{\mathsf{b}_{\mathsf{r}}}
\newcommand{\sy}{\mathsf{s}}
\newcommand{\ev}{\mathsf{e}}
\newcommand{\qq}{\mathsf{q}}
\newcommand{\cons}{\mathsf{cons}}
\newcommand{\consm}{\mathsf{cons}_{\mm}}
\newcommand{\consd}{\mathsf{cons}_{\dd}}
\newcommand{\conss}{\mathsf{cons}_{\sy}}
\newcommand{\conse}{\mathsf{cons}_{\ev}}
\newcommand{\consf}{\mathsf{cons}_{\fw}}

\newcommand{\Spec}{\mathcal{S}}
\newcommand{\Nm}{\mathcal{N}}
\newcommand{\cre}[1]{a^{\dagger}_{#1}}
\newcommand{\ann}[1]{a_{#1}}
\newcommand{\Wt}{\mathsf{W}}

\theoremstyle{definition}
\newtheorem{definition}{Definition}[section]
\newtheorem{requirement}[definition]{Requirement}
\newtheorem{decision}[definition]{Decision}
\newtheorem{obligation}[definition]{Obligation}
\newtheorem{hazard}[definition]{Hazard}
\newtheorem{erratum}[definition]{Erratum}
\newtheorem{remark}[definition]{Remark}
\newtheorem{example}[definition]{Example}
\theoremstyle{plain}
\newtheorem{proposition}[definition]{Proposition}
\newtheorem{lemma}[definition]{Lemma}

\newcommand{\Lft}{\mathsf{L}}
\newcommand{\Rgt}{\mathsf{R}}
\newcommand{\inst}{\mathsf{inst}}
\newcommand{\ERECT}{\mathrm{ERECT}}
\newcommand{\ALLOC}{\mathrm{ALLOC}}
\newcommand{\enc}[1]{\llbracket #1 \rrbracket}
\newcommand{\inp}[3]{\mathsf{for}(\,#1 \leftarrow #2\,)#3}
\newcommand{\out}[2]{#1!(#2)}
\newcommand{\Fam}{\mathcal{F}}

\lstset{
  basicstyle=\ttfamily\footnotesize,
  columns=fullflexible,
  keepspaces=true,
  breaklines=true,
  frame=leftline,
  framesep=6pt,
  xleftmargin=8pt,
  rulecolor=\color{black!30},
  showstringspaces=false,
}

\newcommand{\cstar}{\mathsf{cons}^{*}}

\title{\textbf{\Mat}\\[1mm]
\large The rho combinators as sparse linear algebra:\\
a Rust implementation of the matrix representation,\\
the bulk-synchronous machine, and the transfer operator\\[2mm]
\normalsize Specification v0.5 --- companion to \Comb{} v0.6}
\author{F1R3FLY.io}
\date{23 September 2026}

\begin{document}
\maketitle

\begin{abstract}
\noindent
\Mat{} is the third stage of the route from the rho calculus to a
data-parallel device: it takes a \Comb{} term to the matrix representation of
\cite{matnote} and runs it on a bulk-synchronous machine whose step is a
maximal independent set of redexes. Phase two of the translation erects each
run-time instance at an address, and within the discipline draft~3's
correctness proof needs, that takes a single-premise rule of multiplicity at
least two \cite[Prop.~6.17, Lem.~6.19]{draft3}. The implementation path uses
the context-labelled combinator
$\inst(a,f,C) \mid \mm(a,v) \red \mm(f,\quo{C[v]})$.

\Comb{} v0.6 now agrees \cite[Dec.~5.19]{combspec}: $\inst$ is its default,
and the curried path-indexed constructors are kept behind
\code{curried-erection} as the build \cite{logic} can read. This version
brings \Mat{} into line with it. $\inst$ is the machine's erection with no
feature to enable. Curried artefacts run on the host behind the same
\code{curried-erection} feature, for the cross-scheme obligations \Comb{} has
added \cite[Obl.~11.5, Obl.~11.6]{combspec}. The literal erection remains a
negative control. The context node, $\inst$ and the hole take \Comb's tags,
with every hole collapsed to one reserved name on import, which is sound
because $\inst$ has a single input.

Two points of \Comb{} v0.6's rationale for $\inst$ read differently from the
machine side, and are returned as requests. The depth argument is right as a
count of bulk-synchronous steps --- 2 per instance under $\inst$ against
$\lceil\log_{2} n\rceil + 2$ under curried, and 7 against 10 per unfolding of
the recursion combinator --- but not as the band of an upper-triangular
operator, since every compiled program with an input contains $\ev$ and has
no such operator. And the concern that $\inst$ makes the rule matrix
program-dependent \cite[\S15, item~5]{combspec} does not arise here: this
representation holds the context as a name in the state, so $\inst$ is one
rule for every program. The curried scheme's limited reach (E7) now matters
more than it did, because the curried build is the one the logic reads and
the one \Comb's new obligations run on the whole corpus.
Key words \MUST, \MUSTNOT, \SHOULD, \MAY{} are used in the sense of RFC~2119.
\end{abstract}

\tableofcontents

%======================================================================
\section{Purpose and scope}\label{sec:purpose}

\subsection{The route, and where this document sits on it}

\begin{center}
\begin{tabularx}{\linewidth}{@{}l>{\raggedright\arraybackslash}Xl@{}}
\toprule
Stage & What it does & Specified in\\
\midrule
1 & K0 text to a canonical normal form & \cite{k1ndl1ng}\\
2 & Normal form to combinator term, in two phases & \Comb{} v0.6 \cite{combspec}\\
3 & Combinator term to matrix representation; the bulk-synchronous machine;
    the transfer operator & this document\\
4 & The machine's kernels on a device & Section~\ref{sec:device}, as a
    contract\\
\bottomrule
\end{tabularx}
\end{center}

\noindent
This version aligns the machine with \Comb{} v0.6, which takes $\inst$ as
its default. Otherwise it leaves v0.3's changes to phase two, and the rest of
v0.2 \cite{matv02}, as they were. Where a requirement below says ``as
v0.2'', v0.2's text is normative; where the two differ, this document
governs. Two properties govern everything below. The machine computes exactly one step
per state, seed and configuration, so a host and a device can be compared
trace for trace (\cite[P4]{matnote}); and it runs what stage two emits,
including every instance of every unit, at the width the program has
(\cite[Princ.~10.4]{matnote}).

\subsection{What exists today}

\begin{center}
\begin{tabularx}{\linewidth}{@{}l>{\raggedright\arraybackslash}X@{}}
\toprule
Artefact & What it gives this project\\
\midrule
\code{rho-combinators-draft3} & The calculus and its presentations
(\S3.2--3.3); the constructor family (Def.~3.8, Rem.~3.9); erection at a
computed name (Rem.~3.11); phase one (\S6.2) and phase two (\S6.3, Def.~6.8,
Ex.~6.9); the discipline and the permutation lemma (Def.~6.10, Lem.~6.11);
freshness by size (Lem.~6.14); \emph{new:} the literal erection's failure
(Ex.~6.16), multiplicity (Prop.~6.17, Def.~6.18, Lem.~6.19), the marking
(Rem.~6.20), the deciding term (Ex.~6.21), the two schemes and the
recommendation (\S6.7, Rem.~6.22); the weighted count (Lem.~9.2)\\
\code{f1r3comb-spec-v06} & The term type and shape table; tags, now with
$\inst$, the context node and the hole (Req.~4.11); phase two as a compiler;
the literal erection forbidden (Req.~5.15); the marking analysis
(Req.~5.17); \emph{new:} $\inst$ by default and curried behind
\code{curried-erection} (Dec.~5.19, Rem.~5.22, Rem.~5.23); the context node
(Req.~5.24); filling a context (Req.~5.25); the curried build
(Req.~5.27); both schemes on the corpus, and contexts round-trip
(Obl.~11.5, 11.6)\\
\code{rhocomb-matrix} v1.2 & The state space, incidence, step and operator;
\S14, phase two seen from the machine\\
\code{rhocombmat.py} v1.3, \code{phase2.py}, \code{phase2-results.txt} & The
reference machine with the generated table, $\inst$, and the curried
constructor $\cstar$; the three-arm experiments E5--E7 of
Section~\ref{sec:evidence}. The repository still holds the two-arm versions
(v1.2); the three-arm files accompany this specification and are asked to be
committed (Section~\ref{sec:paper}, item~1)\\
\code{rhocomb-logic} rev.~2 & Components (Def.~2.9); drop elimination
(Lem.~2.8); injective quotation (Prop.~2.10); the weighted count as a formula
(Def.~4.2, Prop.~6.17)\\
\bottomrule
\end{tabularx}
\end{center}

\subsection{Nomenclature}\label{sec:nomen}

\begin{center}
\begin{tabularx}{\linewidth}{@{}l>{\raggedright\arraybackslash}X@{}}
\toprule
Term & Meaning in this document\\
\midrule
\emph{term}, \emph{shape table}, \emph{family} & As in v0.2: a closed
drop-free \code{Term}; \Comb's table of shapes; the constructors a program
uses, $\Fam(P)$\\
\emph{row width}, $w$ & The largest arity in the program's signature\\
\emph{name id}, \emph{row}, \emph{row ref} & As in v0.2\\
\emph{address}, \emph{leaf}, \emph{static channel}, \emph{unit},
\emph{instance} & As in v0.2 and \cite[\S1.3]{combspec}\\
\emph{scoped}, \emph{marked} & A bound name allocated beneath the address; a
name carrying an instance-specific value, in the sense of
\cite[Rem.~6.20]{draft3}\\
\emph{template} & An interned name over the reserved \code{HOLE}: the thing
an erecting rule instantiates at an address\\
\emph{scheme} & How a unit is erected: \emph{curried}
\cite[\S6.7]{draft3}, \emph{$\inst$} \cite[\S6.7]{draft3}, or \emph{literal}
\cite[Def.~6.8]{draft3}\\
\emph{provenance} & The instance a row belongs to, as far as the machine can
tell; metadata, never semantics\\
\emph{mixed} & An atom mentioning scoped names of two addresses\\
\bottomrule
\end{tabularx}
\end{center}

\subsection{Goals}

\begin{enumerate}[leftmargin=1.6em]
\item A total, lossless, round-trippable map from \Comb{} terms to the matrix
representation, for every shape \Comb's shape table generates and every
scheme \Comb{} emits.
\item Redex finding as joins derived from the generated rule table, with three
finders that agree.
\item A bulk-synchronous step defined bit-exactly.
\item Phase-two images erected by the context-labelled combinator $\inst$
run correctly, with \Comb's allocator, freshness checks and marking; curried
artefacts run on the host for \Comb's cross-scheme obligations; the literal
scheme is available as a negative control.
\item The instruments that bear on \Comb's open choices: width with erection
separated, per-scheme cost, cross-instance mixing, and a brute-force oracle
for the marking analysis.
\item On the destructor-free sublattice, the transfer operator, scoped
honestly.
\item \Comb's portability envelope.
\end{enumerate}

\noindent
Non-goals are those of v0.2: device code (a contract only), presentation~B,
weights, the logic, the history unfolding, intern-table collection, and phase
one and two themselves. \Mat{} takes a \code{Term} and \MUSTNOT{} depend on
\code{f1r3comb-lower} or \code{f1r3comb-ir}.

%======================================================================
\section{Position relative to \Comb}\label{sec:position}

\begin{center}
\begin{tabularx}{\linewidth}{@{}l>{\raggedright\arraybackslash}X>{\raggedright\arraybackslash}X@{}}
\toprule
& \Comb{} machine (\code{f1r3comb-run}) & \Mat{} (\code{f1r3comb-par})\\
\midrule
State & \Camp{} store keyed by \code{Hash32} & shape tables of \code{NameId}
rows, width $w$\\
Name identity & content hash of the encoding & structural intern in-run,
content hash at the boundary\\
Event & one comm & one step, a maximal independent set\\
Resolver & store order key; adversarial profile \cite[Obl.~11.2]{combspec}
& priority function; cross-instance profile\\
Rules & generated from the shape table & generated from the same table\\
Erection & $\inst$ by default; curried under \code{curried-erection} &
$\inst$ by default; curried on the host under \code{curried-erection}; literal
as a negative control\\
Addresses & allocator \cite[Req.~9.6]{combspec} & the same allocator\\
\bottomrule
\end{tabularx}
\end{center}

\subsection{Decisions}

\begin{decision}[Presentation A only]\label{dec:aonly}
Unchanged from v0.2. Under B rows are ragged and (G1) of
\cite[\S2.1]{matnote} fails; the A/B comparison \cite[Obl.~11.12]{combspec}
is run on \code{f1r3comb-run}.
\end{decision}

\begin{decision}[A store row holds the quotation of its payload]
\label{dec:qrow}
Unchanged from v0.2: the row for $\qq(a,p)$ is
$(\qq,\mathrm{id}(a),\mathrm{id}(\quo{p}))$, lossless under A, with the shape
tag carrying the $\mm$/$\qq$ distinction, and $\cons_{\qq}$ uniform with every
other member.
\end{decision}

\begin{decision}[Decode by arena read]\label{dec:decode}
Unchanged from v0.2, and still at odds with \cite[Req.~9.4]{combspec}
(Section~\ref{sec:paper}, item~10).
\end{decision}

\begin{decision}[The rule table is generated from \Comb's shape table]
\label{dec:ruletable}
Unchanged from v0.2. The generator lives in \code{f1r3comb-term}; \Mat{} lists
no rules of its own.
\end{decision}

\begin{decision}[$\inst$ is the erection; curried behind the same feature as \Comb]
\label{dec:schemes}
The implementation path erects by the context-labelled combinator $\inst$,
and so does \Comb{} v0.6 by default \cite[Dec.~5.19]{combspec}. \Mat{}
\MUST{} support $\inst$ unconditionally. Every measurement that bears on the
path --- width, cost per instance, conformance, the device conformance
corpus --- \MUST{} be made on $\inst$ artefacts, and the device contract
targets $\inst$ alone (Requirement~\ref{req:device}). \Mat{} runs whatever
scheme the artefact header names, and under \code{curried-erection}, the
feature name \Comb{} uses, \MUST{} also run curried artefacts on the host.
That support exists for \Comb's cross-scheme obligations
\cite[Obl.~11.5, Obl.~11.6]{combspec} and for the logic tooling's build
\cite[Req.~5.27]{combspec}; it is not a second execution path. It \MUST{} support the
literal scheme under a test-only feature \code{literal-erection}, which no
product crate enables, because \cite[Obl.~11.3]{combspec} requires the broken
construction as a negative control and \Comb{} is forbidden to emit it
\cite[Req.~5.15]{combspec}. It \MUST{} accept units that \Comb's marking
analysis has compiled with the fixed-arity family \cite[Req.~5.17]{combspec},
reading the marking from the header (Section~\ref{sec:marking}).
\end{decision}

\begin{decision}[One kernel, whose primary client is $\inst$; contexts are data]
\label{dec:onekernel}
A curried constructor $\cstar_{A}[\vec w](t,f)$ and an $\inst(a,f,C)$ both
consume one message and deliver the quotation of a template instantiated at
its payload. \Mat{} \MUST{} implement both by one rule shape with one encode,
\code{Instantiate\{template, with\}}, and distinguish them by a validity class
on the template (Requirement~\ref{req:template}). \Comb{} v0.6 records that
$\inst$ adds one term operation to its \Camp{} machine
\cite[Rem.~9.2]{combspec}; on this machine there is nothing to add, since the
curried scheme's leaves are built by the same substitution and both schemes
are the same kernel.

Because the template is the third column of the row, a context is a name in
the state, not part of the rule. $\inst$ is therefore one rule for every
program, and the rule table is independent of the program's units; so is
$\cstar$'s. This answers the second half of \cite[\S15, item~5]{combspec} for
this representation (Section~\ref{sec:paper}, item~4).
\end{decision}

%======================================================================
\section{Phase two, with the evidence}\label{sec:evidence}

\subsection{The obstruction, now in the paper}

\cite[\S6.6]{draft3} shows the literal erection fails the discipline on a
single instance (Ex.~6.16), that the obvious repair regresses, and that the
obstruction is structural: under the discipline extended to all rules every
reachable atom holds at most one occurrence of an address
(Prop.~6.17), so erecting any atom with two scoped names needs a rule of
multiplicity at least two (Lem.~6.19). v0.2's Section~3.2 is discharged by
this and \Comb{} v0.5 withdraws its two-stage construction
\cite[Rem.~5.16]{combspec}. \cite[\S6.7]{draft3} supplies two single-premise
rules of the required multiplicity:
\[
  \cstar_{A}[\vec w](t,f) \mid \mm(t,\sigma) \red
  \mm(f,\quo{A(\sigma\!\cdot\!w_{1},\dots,\sigma\!\cdot\!w_{n})}),
  \qquad
  \inst(a,f,C) \mid \mm(a,v) \red \mm(f,\quo{C[v]}).
\]

\subsection{Three arms}

\cite[Rem.~5.20]{combspec} still records that the committed measurements
compare the literal scheme with $\inst$ and have no curried arm, and that a
third arm would let the figures be compared directly rather than argued. The
three-arm \code{phase2.py} that accompanies this specification runs it, with
the curried scheme implemented as \cite[Rem.~5.21]{combspec} describes it: no allocation gadget, one $\dd$ per erected atom fanning the
address out on static channels, one $\cstar$ and one $\ev$ per atom. Its
templates allow a static name at an argument position, as
\cite[Ex.~6.16]{draft3}'s $\dd(a,p_{0},q_{1})$ requires.

\begin{center}\small
\begin{tabular}{@{}lrrrr@{}}
\toprule
\multicolumn{5}{@{}l}{\textbf{E5.} Unit $\{\fw(z_{1},z_{2}),\ \kk(z_{1})\}$;
runs with a mixed atom, of 50}\\
\midrule
instances & 2 & 4 & 8 & 16\\
\midrule
literal, maximal progress / sequential & 29 / 16 & 48 / 38 & 50 / 49 & 50 / 50\\
curried, either schedule & 0 & 0 & 0 & 0\\
$\inst$, either schedule & 0 & 0 & 0 & 0\\
\midrule
parallel steps to erect: literal 5, curried 3, $\inst$ 2\\
\bottomrule
\end{tabular}
\end{center}

\begin{center}\small
\begin{tabular}{@{}lrrr@{}}
\toprule
\multicolumn{4}{@{}l}{\textbf{E6.} $D_{x}$ of \cite[Ex.~6.9]{draft3}, 200
unfoldings, per unfolding}\\
\midrule
scheme & parallel steps & sequential steps & new names\\
\midrule
curried & 10 & 32.0 (29--35) & 22\\
$\inst$ & 7 & 10.0 (9--11) & 16\\
\bottomrule
\end{tabular}
\end{center}

\begin{center}\small
\begin{tabular}{@{}lrr@{}}
\toprule
\multicolumn{3}{@{}l}{\textbf{E7.} Store $\qq(b,\quo{(\fw(c_{1},c_{3}) \mid
\kk(c_{2}))})$, $b$ static, $c_{i}$ scoped; mixed stores}\\
\midrule
scheme & 2 instances & 8 instances\\
\midrule
curried, payload joined by $\cons_{\mid}$ & 44 / 100 & 347 / 400\\
$\inst$ & 0 / 100 & 0 / 400\\
\bottomrule
\end{tabular}
\end{center}

Three conclusions, each a requirement below.

\emph{Both conforming schemes are correct for top-level atoms.} E5 answers
\cite[Rem.~5.20]{combspec}'s question on conformance: no mixing under either,
at any count, under either schedule. Conformance does not separate them;
depth, cost and reach do.

\emph{$\inst$ erects in fewer bulk-synchronous steps.} Two per instance,
against $\lceil\log_{2} n\rceil + 2$ for a curried unit of $n$ atoms (the
duplicator tree, then $\cstar$, then $\ev$); 7 against 10 per unfolding of
$D_{x}$. This is \cite[Rem.~5.22]{combspec}'s depth argument, measured in the
unit it holds in (Section~\ref{sec:paper}, item~2).

The three together are why the implementation path uses $\inst$.

\emph{On this machine $\inst$ also does less work.}
\cite[Rem.~5.22]{combspec} holds that the work figure does not change between
the schemes: $\inst$ spends $O(|\mathrm{unit}|)$ in one step, recomputing
encodings up the rebuilt term, and curried spends the same across its steps.
Structural interning computes no hashes (Requirement~\ref{req:hashfree}), and
the substitution visits only names that mention the hole
(Requirement~\ref{req:subst}), so the work here is names interned: 16 per
unfolding under $\inst$ against 22 under curried, where each $\cstar$ builds
its leaves by $d$ interns from the address. Both schemes are charged by names
built, as \cite[Req.~5.25]{combspec} requires, and on that measure $\inst$ is
cheaper here.

\emph{The curried scheme does not reach every store.} Its template places the
address at argument positions of one atom. A gate store whose body mentions
scoped names in two atoms --- which is the body of any input whose bound name
occurs in two atoms of its continuation --- needs its payload assembled from
two address-bearing pieces, and the only rule that assembles them is a join.
The store atom then mentions two marked names, which is
\cite[Rem.~6.20]{draft3}'s own criterion, so Lemma~6.19 applies to it; the
multiplicity it needs is beneath a quotation, which the curried rule does not
supply. \cite[Req.~4.9]{combspec}'s recipe (build the payload by $\cons_{\ev}$,
feed it to $\cons_{\qq}$) conforms exactly when the store's subject is static
and the payload has one scoped atom, which is the case of
\cite[Ex.~6.9]{draft3}; E6 uses it there. E7 is the case beyond it.

\begin{hazard}[Curried stores]\label{haz:curriedstores}
Under \code{curried-erection}, a unit whose gate store's payload mentions
scoped names in more than one atom, or beneath a further quotation, cannot be
erected within the discipline, and a compiler that falls back to a join
produces the mixing E7 measures. This matters more under \Comb{} v0.6 than
before: the curried build is the one the logic reads
\cite[Req.~5.27]{combspec}, and \cite[Obl.~11.5, Obl.~11.6]{combspec} run it
on the whole corpus and compare it with $\inst$, which such units will fail
for reasons of the scheme, not the compiler. Whether that mixing is observable is the deciding-term
question of \cite[Ex.~6.21]{draft3} applied to stores, which the marking
analysis does not yet ask (Section~\ref{sec:paper}, item~7). \Mat{} \MUST{}
detect such stores at deploy under the curried scheme and report them as
\code{comb-mat-curried-store}, a warning, naming the unit.
\end{hazard}

\subsection{Templates}

\begin{requirement}[One representation for both schemes]\label{req:template}
A template \MUST{} be an interned name over the reserved \code{NameId::HOLE}.
Both erecting rules \MUST{} be rows of arity three --- $(t,f,T)$ for $\cstar$,
$(a,f,C)$ for $\inst$ --- with the template in the third column, so the row
width is four under either scheme. A $\cstar$ template \MUST{} satisfy the
curried validity class: it is a single atom, and each argument is either a
path beneath \code{HOLE} (a $\Lft/\Rgt$ spine) or a name that does not mention
\code{HOLE}. The check runs at import and at every firing in debug builds;
a violation is \code{comb-mat-template}. An $\inst$ template is any name over
\code{HOLE}.
\end{requirement}

\begin{requirement}[Contexts and holes on import and export]\label{req:holeidx}
\Comb's context node (tag \code{0x61}) and hole (tag \code{0x62}, carrying a
canonical index) \cite[Req.~4.11, Req.~5.24]{combspec} \MUST{} be imported as
a template over the single reserved \code{HOLE}, every hole index collapsed.
This is sound because $\inst$ has one input, so every hole is filled with the
same name, and it is necessary for structural interning: two occurrences of
the same leaf beneath distinct hole indices would otherwise be distinct names,
and a template would share nothing. Export \MUST{} renumber holes in encoding
order, so that export--import and import--export round-trip against \Comb's
encoding (Section~\ref{sec:paper}, item~6).
\end{requirement}

\begin{requirement}[The hole]\label{req:hole}
As v0.2: \code{HOLE} has no counterpart in the term type outside a template
position, is rejected by import anywhere else (\code{comb-mat-hole}), and no
rule's conclusion may contain it. Its encoding is \Comb's hole tag
\code{0x62} (Requirement~\ref{req:holeidx}); $\inst$ is \code{0x60}.
$\cstar$ has no tag yet (Section~\ref{sec:paper}, item~7).
\end{requirement}

\begin{requirement}[Instantiation]\label{req:subst}
The encode \MUST{} compute $T\{v/\code{HOLE}\}$ over the interned form,
hereditarily beneath quotation, with a per-firing memo keyed by name id. The
set of names in $T$ that mention \code{HOLE} \MUST{} be precomputed once per
template, at import, and the traversal \MUST{} visit only those, reusing every
other name of $T$ unchanged; this is \cite[Req.~5.25]{combspec}'s
precomputed hole positions in interned form, and its cost is the number of
hole-bearing names of $T$, not the size of the unit. Result names are interned in
canonical order (Requirement~\ref{req:insertorder}). The report \MUST{}
attribute the names built to the scheme, which is the work figure of E6.
\end{requirement}

\begin{requirement}[The measured costs are regressions, $\inst$ first]
\label{req:dxlaw}
\code{f1r3comb-check} \MUST{} reproduce E6 as laws over the number of
unfoldings to $10^{4}$: 7 parallel steps and 16 names per unfolding under
$\inst$, with a sequential mean of 10, as the primary regression, matching
\cite[Obl.~11.3]{combspec}; and 10, 22 and 32 under \code{curried-erection}. A golden
file \MUSTNOT{} be compared across schemes \cite[Dec.~5.19]{combspec}.
\end{requirement}

\subsection{Marking, and when a crossing is sanctioned}\label{sec:marking}

\cite[Req.~5.17]{combspec} compiles a unit with the fixed-arity family when no
atom of its phase-one image mentions two marked names, since crossings there
permute which instance handles which value and every answer is right. On such
a unit this machine will observe crossings --- that is what the analysis
licenses --- and must not report them as defects.

\begin{requirement}[The header's marking is read, and checked]
\label{req:marking}
\Mat{} \MUST{} read the marking from the artefact header and carry, for each
scoped name, whether it is marked. A row mentioning scoped names of two
addresses is \emph{mixed}; it is \emph{sanctioned} if at most one of those
names is marked and \emph{unsanctioned} otherwise. Unsanctioned mixed rows
\MUST{} raise a trace event at every trace level and be counted in the run
report; under a conforming scheme with a correct marking there are none, and
the test suite \MUST{} assert it under \code{CrossInstance}.
\end{requirement}

\begin{requirement}[The oracle for the marking analysis]\label{req:oracle}
\cite[Obl.~11.4]{combspec} checks the marking against exhaustive scheduling
at two and three instances. \code{f1r3comb-op}'s bounded exploration
(Section~\ref{sec:op}) \MUST{} serve as that oracle: given a unit, an instance
count and a depth bound, it explores every interleaving, collects the encoded
observations at quiescent markings, and reports whether they differ. It
\MUST{} classify $R = \inp{y}{a}{\out{y}{\drp{y}}}$ as unsafe and
$\inp{y}{a}{(\drp{y} \mid \drp{y})}$ as safe \cite[Req.~5.18]{combspec},
and \SHOULD{} also be run on E7's store, which neither \Comb{} nor the draft
has yet classified.
\end{requirement}

\subsection{Provenance}\label{sec:prov}

\begin{definition}[Provenance]\label{def:prov}
As v0.2: every row carries \code{Prov}, \code{STATIC} on import; a produced
row takes its premises' common non-static tag, \code{STATIC} if all are
static, \code{MIXED} otherwise; and a row naming a leaf takes that leaf's
root's tag from the intern table's \code{root} column. For a template
instantiation the tag is the consumed address's.
\end{definition}

\begin{requirement}[Provenance is metadata]\label{req:provmeta}
As v0.2: it \MUSTNOT{} affect redexes, selection under
\code{MaximalProgress}, or conclusions; it appears in traces as the root
address's content hash. \code{MIXED} rows are classified sanctioned or
unsanctioned by Requirement~\ref{req:marking}.
\end{requirement}

%======================================================================
\section{The representation}\label{sec:rep}

v0.2's representation is unchanged: structural interning with no hash on the
hot path; \code{size}, \code{root} and \code{depth} columns; ids never leaving
the process; the budget, including address depth; structure-of-arrays shape
tables with a provenance column; import, export and round trip; the
canonical \code{.cmat}. Three points are restated.

\begin{requirement}[Structural interning]\label{req:hashfree}
As v0.2, with the observation of E6 added: prefix sharing is by construction,
and per-unfolding name cost is constant under both schemes (16 and 22).
\end{requirement}

\begin{requirement}[Width]\label{req:soa}
$w$ is read from the header and instantiated for $w \in \{4,5,6\}$. Under the
curried scheme and under $\inst$, $w = 4$: base shapes have arity at most
three, $\cons_{\dd}$ and $\cons_{\sy}$ four, the erecting rows three. $w = 5$
arises only under \code{literal-erection} with an (o5) occurrence, which is a
negative control.
\end{requirement}

\begin{requirement}[\code{.cmat}]\label{req:cmat}
As v0.2, with the header extended by the scheme, the marking, and the
template names; \code{HOLE} relabelled to index $0$.
\end{requirement}

%======================================================================
\section{The rule table}\label{sec:rules}

\begin{requirement}[What the generator emits]\label{req:gen}
As v0.2 --- the seven routing and destructor rules, $\mathsf{quote}$,
$\mathsf{make}_{\mid}$, one rule per member of $\Fam(P)$ --- and, per the
header's scheme, the erecting rule: under curried, $\cstar$ with one loud
premise at slot 0 and conclusion
\code{Encode\{deliver: Slot(1), node: Instantiate\{template: Slot(2), with:
Payload(0)\}\}}; under $\inst$ the same with its own shape; under literal,
nothing further, the family doing the erecting. Classes, join plans, the
grading and $w$ are derived from the emitted table.
\end{requirement}

\begin{requirement}[Only a forwarder reads a store]\label{req:quietjoin}
As v0.2, a property of the generator.
\end{requirement}

\begin{requirement}[Classes and plans are checked]\label{req:plans}
As v0.2, per scheme: one class-(c) rule; one class-(b) rule per member of
$\Fam(P)$, plus $\mathsf{make}_{\mid}$ and the scheme's erecting rule; seven
class-(a) rules.
\end{requirement}

\begin{requirement}[Grading]\label{req:grading}
As v0.2. The erecting rules, like the family, have weight one.
\end{requirement}

%======================================================================
\section{Redex finding, selection, firing}\label{sec:machine}

The sparse layer, the three finders, the canonical key, the distinctness and
symmetry masks, $n$-way joins over the generated signature, the conflict
hypergraph, the \code{mix64} priority, the greedy maximal independent set and
its equality with parallel rounds, the four resolvers including
\code{CrossInstance}, race visibility, the order of effects, and debug-build
soundness checking are unchanged from v0.2.

\begin{requirement}[Order of effects]\label{req:insertorder}
As v0.2: sort the step by canonical key; compute conclusions in that order,
interning new names --- including those an instantiation creates, in a fixed
traversal of the template --- in that order; remove consumed rows by stable
compaction; append produced rows in the order computed. Parallel backends
\MUST{} assign ids as if sequentially.
\end{requirement}

\begin{hazard}[Contention at statics, per scheme]\label{haz:statics}
Under $\inst$, $n$ instances contend at one static channel $r$ per unit with a
one-premise rule: $n^{2}$ redexes, a perfect matching in one step. Under
curried the contention is spread over the duplicator tree and the $\cstar$
channels, each one-premise, so each is a perfect matching; the enumerated
count is $n^{2}$ per erected atom. Under literal the multi-premise
constructors enumerate $n^{m}$. The enumerated-to-fired ratio \MUST{} be
reported per static channel and per scheme.
\end{hazard}

%======================================================================
\section{Addresses and freshness}\label{sec:addr}

\begin{requirement}[One allocator for both machines]\label{req:oneallocator}
As v0.2: \Comb's allocator \cite[Req.~9.6]{combspec}, factored into
\code{f1r3comb-term}.
\end{requirement}

\begin{requirement}[Deploy checks]\label{req:deploychecks}
As v0.2, citing \Comb{} v0.6: root address larger than every name of the image
(\cite[Req.~5.28]{combspec}), static channels likewise
(\cite[Req.~5.31]{combspec}), root addresses prefix-incomparable
(\cite[Req.~5.32]{combspec}), with \Comb's diagnostics. Under $\inst$ the static channels are each unit's $r$ and $f$; under
\code{curried-erection} they also include every channel of the duplicator
tree.
\end{requirement}

%======================================================================
\section{The operator layer}\label{sec:op}

v0.2's operator layer is unchanged in content --- species, markings,
successors with mass-action coefficients, exploration, the materialised
$T$ graded by $\Wt$, closure, $T^{\top}\chi$ --- and so is its scope: every
compiled unit with an input contains $\ev$, so triangularity and closure apply
to none of them, which \Comb{} now reports too \cite[Req.~10.2]{combspec}. One
use is added.

\begin{requirement}[Bounded multi-instance exploration]\label{req:multiexplore}
\code{explore} \MUST{} accept an instance count $k \in \{2,3\}$ and a depth
bound, deploy $k$ instances at prefix-incomparable addresses, and enumerate
every interleaving to the bound, returning the set of encoded observations at
quiescent markings. This is Requirement~\ref{req:oracle}'s oracle. The
exploration is truncated by construction and \MUST{} say so; it makes no
triangularity or closure claim.
\end{requirement}

%======================================================================
\section{Instrumentation}\label{sec:instr}

\begin{requirement}[Reports]\label{req:reports}
As v0.2, with erection attributed per scheme --- the $\cstar$ or $\inst$
firings and their $\ev$, and under curried the duplicator tree --- and with
names built per instance, sanctioned and unsanctioned mixed rows, and curried
stores detected (Hazard~\ref{haz:curriedstores}).
\end{requirement}

\begin{requirement}[The width measurement]\label{req:widthcmd}
\code{f1r3x comb-width} \MUST{} report, per program, the width distribution
under $\inst$ with erection separated from computation, and \MAY{} report the
curried scheme alongside for comparison. Under $\inst$
erection is two steps per instance; under curried it is the duplicator tree's
depth plus two. Whether either adds width without depth on real programs is
the measurement \cite[Obl.~13.7]{matnote} asks for.
\end{requirement}

\begin{requirement}[Subcommands]\label{req:fx}
As v0.2, with \code{-{}-scheme} on \code{comb-run --machine=matrix} and
\code{comb-width}, and \code{-{}-instances} on \code{comb-op}.
\end{requirement}

%======================================================================
\section{Architecture}\label{sec:arch}

The crates are v0.2's: \code{f1r3comb-term} (\Comb's, plus the rule generator
and the allocator), \code{f1r3comb-mat}, \code{f1r3comb-par},
\code{f1r3comb-op}, and the shared \code{f1r3comb-check}; none depends on
\code{f1r3comb-ir}, \code{f1r3comb-lower} or \code{f1r3comb-run}. The
test-only \code{literal-erection} feature \MUST{} live in
\code{f1r3comb-check}'s emitter and \code{f1r3comb-par}'s dev-dependencies
only.

\subsection{The device seam}\label{sec:device}

\begin{requirement}[Device kernels]\label{req:device}
As v0.2, with the erection fixed: a device backend \MUST{} target $w = 4$ and
the $\inst$ erection. Its substitution kernel is a gather over the template's
interned form with a batched insert-or-get of the results; its one erecting
join is a one-premise match at each unit's static $r$. Curried artefacts
\MUSTNOT{} reach a device --- they add a duplicator tree per unit, build more
names (E6) and cannot erect every store (E7) --- and neither may
\code{literal-erection}. A curried artefact submitted to a device backend
\MUST{} be rejected with \code{comb-mat-scheme}. It is accepted only on
identical step traces to \code{HostSequential} on the $\inst$ conformance
corpus, and no work on it \MAY{} begin before Obligation~\ref{obl:width} is
discharged. v0.2's gate on phase two's primitive is discharged by
\cite[\S6.7]{draft3}: the primitive is $\inst$.
\end{requirement}

%======================================================================
\section{Correctness obligations and testing}\label{sec:test}

\begin{obligation}[Carried from v0.2]\label{obl:carried}
Round trip; finder agreement over every generated shape, now including
$\cstar$; Python parity on E1--E4; trace certification; operator correctness;
determinism across backends, core counts and WebAssembly; the invariants of
Requirement~\ref{req:grading}; device conformance.
\end{obligation}

\begin{obligation}[The phase-two experiments, three arms]\label{obl:phase2}
\code{f1r3comb-check} \MUST{} reproduce \code{phase2-results.txt}: the E5
mixing rates for the literal scheme within a tolerance, and exactly zero for
curried and $\inst$ at every instance count under \code{MaximalProgress},
sequential choice and \code{CrossInstance}; the E6 laws exactly
(Requirement~\ref{req:dxlaw}); and E7, zero for $\inst$ and a nonzero rate
for the curried join. This is \cite[Obl.~11.3]{combspec} with the curried arm
it lacked, and the literal scheme is its negative control.
\end{obligation}

\begin{obligation}[The $W$ regression and the deciding terms]\label{obl:W}
$W$ \cite[Obl.~11.1]{combspec}, its three-instance variant, and both halves of
\cite[Req.~5.18]{combspec}, run under \code{MaximalProgress} and
\code{CrossInstance} for 100 seeds each, under each conforming scheme: the
observation at $\enc{\quo{X}}$ \MUST{} match \Bell's on every seed and no row
\MUST{} be unsanctioned-mixed. Under \code{literal-erection} the rate of the
crossed residue is recorded as a negative control.
\end{obligation}

\begin{obligation}[Both schemes, and contexts round-trip]\label{obl:schemes}
\cite[Obl.~11.5, Obl.~11.6]{combspec} on this machine: compile the corpus
under $\inst$ and under \code{curried-erection}, run both, and compare
observations; and for every unit, compare the atoms $\inst$ erects with the
atoms the curried erection produces when run to completion. Units in the
class of Hazard~\ref{haz:curriedstores} \MUST{} be reported separately, since
the curried build cannot erect them within the discipline and a divergence
there is the scheme's, not the implementation's.
\end{obligation}

\begin{obligation}[The marking oracle]\label{obl:marking}
Requirement~\ref{req:oracle} run on \Comb's marking corpus; any disagreement
with \Comb's analysis \MUST{} be reported to \Comb{} as a counterexample
\cite[Obl.~11.4]{combspec}.
\end{obligation}

\begin{obligation}[Addresses and differential testing]\label{obl:addr}
As v0.2, citing \cite[Obl.~11.9, 11.11, 11.14]{combspec}. \Comb{} now emits
both conforming schemes, so v0.2's test-only $\inst$ emitter is retired.
\end{obligation}

\begin{obligation}[The width of real programs]\label{obl:width}
\cite[Obl.~13.7]{matnote}, on $\inst$ artefacts, with erection separated.
Gates the device seam.
\end{obligation}

%======================================================================
\section{Errors and diagnostics}\label{sec:errors}

As v0.2, with three additions:

\begin{center}
\begin{tabularx}{\linewidth}{@{}ll>{\raggedright\arraybackslash}X@{}}
\toprule
Code & Where & Meaning\\
\midrule
\code{comb-mat-template} & import, step & A $\cstar$ template outside the
curried validity class (Requirement~\ref{req:template})\\
\code{comb-mat-curried-store} & deploy & Warning: a curried unit whose store
payload needs a join (Hazard~\ref{haz:curriedstores})\\
\code{comb-mat-scheme} & deploy & A scheme the backend does not carry: curried
or literal on a device (Requirement~\ref{req:device})\\
\bottomrule
\end{tabularx}
\end{center}

%======================================================================
\section{Performance, portability, determinism}\label{sec:perf}

As v0.2, with the targets stated for $\inst$: $10^{4}$ unfoldings of $D_{x}$
with no rise in per-unfolding step time; 64 concurrent instances of a
$10^{3}$-atom unit erected in two steps. For comparison only, the curried
scheme would take $\lceil\log_{2}(\text{atoms})\rceil + 2$.

%======================================================================
\section{What the implementation asks of the papers}\label{sec:paper}

\paragraph{Discharged since v0.4.} \Comb{} v0.6 makes $\inst$ its default
\cite[Dec.~5.19]{combspec}, which was v0.4's first request; its mixing
obligation now asserts $\inst$'s figures \cite[Obl.~11.3]{combspec}, which was
the second; $\inst$, the context node and the hole have tags
\cite[Req.~4.11]{combspec}; and the question of which scheme the logic reads
is now stated by \Comb{} in its sharpest form \cite[\S14, item~5]{combspec}.
Earlier discharges are as recorded in v0.3 and v0.4.

\paragraph{Open.}
\begin{enumerate}[leftmargin=1.8em]

\item \textbf{Commit the three-arm measurements.} \cite[Rem.~5.20]{combspec}
and \cite[\S14, item~2]{combspec} still say the curried arm does not exist.
It does: \code{rhocombmat.py} v1.3, \code{phase2.py} and
\code{phase2-results.txt} accompanying this specification run E5--E7.
Proposed: commit them, and replace Remark~5.20 with the figures of
Section~\ref{sec:evidence}.

\item \textbf{\Comb{} v0.6, Remark~5.22, depth.} The argument that decides
\Comb's choice is stated as keeping ``the band of the upper-triangular matrix
shallow and constant'', with the curried scheme costing $n$ steps for an
$n$-atom unit. Two corrections, neither of which weakens the choice. Every
compiled program with an input contains $\ev$, so it has no upper-triangular
transfer operator (\cite[Thm.~6.1(ii), \S14.2]{matnote}), which \Comb{}
itself records \cite[Req.~10.2]{combspec}; the depth that matters is the
number of bulk-synchronous steps. And in that unit the curried scheme costs
$\lceil\log_{2} n\rceil + 2$ steps, not $n$, because its $n$ constructors fire
in parallel after the duplicator tree (E5: 3 against 2; E6: 10 against 7).
Proposed: restate the argument in steps, with the measured figures.

\item \textbf{\Comb{} v0.6, Remark~5.22, ``a fixed linear map''.} Right in
the sense that matters: every rule is a ladder-operator monomial schematic in
names, and $C[v]$ is a class-(b) encode like any constructor's, a function
on name indices (\cite[\S5.1]{matnote}). Proposed: say so in those terms,
since ``fixed linear map on the vector encoding'' suggests the state vector,
which $C[v]$ does not act on directly.

\item \textbf{\Comb{} v0.6, Remark~5.22, ``one rule row'', and open
question~5.} \Comb{} argues that $\inst$ is one rule row where the curried
family needs a row per shape and path tuple, and then notes that one row per
context makes the rule matrix program-dependent after all. In this
representation neither holds: the template is the third column of the row, a
name in the state, so $\inst$ is one rule for every program and $\cstar$ is
one rule too (Decision~\ref{dec:onekernel}). The rule table is
program-independent under both schemes. Proposed: close open question~5's
second half on that basis, and rest the fourth argument on the step count
rather than the rule count.

\item \textbf{\Comb{} v0.6, Remark~5.22, ``what does not change''.} On this
machine the work figure does change: $\inst$ builds 16 names per unfolding of
$D_{x}$ against curried's 22, because structural interning hashes nothing and
the substitution visits only hole-bearing names (Requirement~\ref{req:subst}).
The requirement to charge by size \cite[Req.~5.25]{combspec} stands and is
met by reporting names built.

\item \textbf{Hole indices} \cite[Req.~5.24]{combspec}. The hole tag carries
a canonical index, but $\inst$ has one input, so every hole is filled with
the same name and the index distinguishes nothing. It does cost sharing: two
occurrences of one leaf beneath different indices encode, hash and intern as
different names. \Mat{} collapses the indices on import
(Requirement~\ref{req:holeidx}). Proposed: an unindexed hole, unless a
multi-input context is intended, in which case say so.

\item \textbf{The curried scheme's reach}, now load-bearing
(Hazard~\ref{haz:curriedstores}). Under v0.6 the curried build is the logic
tooling's target \cite[Req.~5.27]{combspec} and
\cite[Obl.~11.5, Obl.~11.6]{combspec} run it on the whole corpus. Units whose
gate stores mention scoped names in two atoms cannot be erected by it within
the discipline (E7), so those obligations will fail on them for reasons of the
scheme, and \cite[Obl.~11.6]{combspec}'s round trip is defined only where the
curried erection reaches. Proposed: state the limit in draft~3's \S6.7 and
\Comb's Requirement~5.27, exclude or flag such units in both obligations, and
count it in the decision of \cite[\S14, item~5]{combspec}: a logic read over
the curried build does not see those programs at all. $\cstar$ also has no
tag in \cite[Req.~4.11]{combspec}.

\item \textbf{\Comb{} v0.6, Requirement~9.6} still describes the root
address as ``of a shape outside that program's family''; freshness is by size
(Req.~5.28, \cite[Lem.~6.14]{draft3}). Carried.

\item \textbf{\Comb{} v0.6, Requirement~10.4} still cites \cite{logic}'s
weighted count as ``Def.~6.6''; it is Def.~4.2 and Prop.~6.17. Carried.

\item \textbf{\Comb{} v0.6, Requirement~9.4} (decode by projection): scope it
to builds where \code{presentation-b} is reachable. Carried.

\item \textbf{The allocator in \code{f1r3comb-term}}, and \textbf{a Merkle
form of the content hash} for $O(1)$ hashing of new addresses on
\code{f1r3comb-run}. Carried.

\end{enumerate}

%======================================================================
\section{Open questions}

\begin{enumerate}[leftmargin=1.8em]
\item \textbf{The logic and $\inst$} \cite[\S14, item~5]{combspec}: a
structural production for the erecting row, or a logic of the curried build
that does not see E7's class of programs.
\item \textbf{Width under $\inst$} on real programs, with erection separated:
does two-step erection add width without depth, as E5 and E6 predict?
\item \textbf{Cost accounting for $\inst$}: charging by names built rather
than by step, as \cite[Rem.~5.22]{combspec} anticipates.
\item \textbf{Resolver fidelity and fairness} (\cite[Obl.~13.1,
Obl.~13.6]{matnote}).
\item \textbf{Address compaction} \cite[\S15, item~1]{combspec}.
\item \textbf{Garbage}, \textbf{the history unfolding}, \textbf{weights}: as
v0.2.
\end{enumerate}

%======================================================================
\section{Delivery}

\begin{center}
\begin{tabularx}{\linewidth}{@{}ll>{\raggedright\arraybackslash}X@{}}
\toprule
M & Deliverable & Done when\\
\midrule
X0 & \code{f1r3comb-mat} & As v0.2, with templates and \code{HOLE}\\
X1 & Finders & As v0.2, over the generated signature including the erecting
rows\\
X2 & \code{f1r3comb-par}, sequential & Resolvers; the instantiation kernel,
$\inst$ first, then the curried validity check; deploy checks; marking read
from the header; the E6 laws, $\inst$ primary\\
X3 & Phase-two evidence & The three arms of Obligation~\ref{obl:phase2}; $W$
and the deciding terms under $\inst$, and under curried for comparison;
reported to \Comb{} with the requests of Section~\ref{sec:paper}, items~1 and~7\\
X4 & The marking oracle & Multi-instance bounded exploration
(Requirement~\ref{req:multiexplore}); Obligation~\ref{obl:marking} on
\Comb's corpus\\
X5 & The width gate & Under $\inst$, erection separated; needs \Comb{} M2.
\textbf{Gate} on X8\\
X6 & Threaded host & Determinism across cores and \code{wasm32}\\
X7 & \code{f1r3comb-op} & As v0.2, scoped\\
X8 & Device conformance & $\inst$ only; external backend only after
Obligation~\ref{obl:width}\\
\bottomrule
\end{tabularx}
\end{center}

%======================================================================
\section{Changes from v0.4}\label{sec:v4delta}

\begin{center}
\begin{tabularx}{\linewidth}{@{}l>{\raggedright\arraybackslash}X@{}}
\toprule
v0.4 & v0.5\\
\midrule
Companion to \Comb{} v0.5, which defaulted to curried & Companion to \Comb{}
v0.6, which defaults to $\inst$; citations renumbered\\
$\inst$ behind \code{inst-erection}, enabled by default & $\inst$
unconditional; curried behind \code{curried-erection}, \Comb's name\\
Curried support for compatibility, to be retired & Curried support for
\Comb's cross-scheme obligations and the logic tooling's build\\
\code{HOLE} with a tag to be allocated & \Comb's tags for $\inst$, context
and hole; hole indices collapsed on import, renumbered on export
(Requirement~\ref{req:holeidx})\\
Substitution over every name of the template & Over hole-bearing names only,
precomputed per template\\
--- & Contexts are data: $\inst$ and $\cstar$ each one rule for every program
(Decision~\ref{dec:onekernel})\\
--- & Obligation~\ref{obl:schemes}: both schemes, and contexts round-trip, on
this machine\\
Eleven asks, the first to make $\inst$ \Comb's default & Discharged; eleven
open, six new, on v0.6's rationale, hole indices, and the curried reach\\
\bottomrule
\end{tabularx}
\end{center}

\noindent
v0.3's and v0.4's changes carry forward.

%======================================================================
\begin{thebibliography}{99}

\bibitem{matnote} L.\,G. Meredith.
\emph{Rho combinators as sparse linear algebra.} Version~1.2,
\code{publications/rho-combinators/rhocomb-matrix}, with \code{rhocombmat.py},
\code{experiments.py} and \code{phase2.py}, September 2026.

\bibitem{matv02} F1R3FLY.io.
\emph{F1R3Comb-Mat: the rho combinators as sparse linear algebra.}
Specification v0.2, \code{publications/rho-combinators/f1r3comb-mat-spec-v02},
23 September 2026.

\bibitem{combspec} F1R3FLY.io.
\emph{F1R3Comb: a compiler from kernel rho to the name-free rho combinators,
and a machine that runs the result.} Specification v0.6,
\code{publications/rho-combinators/f1r3comb-spec-v06}, 23 September 2026.

\bibitem{draft3} L.\,G. Meredith and M. Stay.
\emph{Name-free combinators for the rho calculus: a repair, a linear encoding,
and the one operation that was missing.} Draft~3,
\code{publications/rho-combinators}, September 2026.

\bibitem{logic} F1R3FLY.io.
\emph{A spatial-behavioural logic for the rho combinators.} Draft~2,
revision~2, \code{publications/rho-combinators/rhocomb-logic}, 23 September
2026.

\bibitem{k1ndl1ng} F1R3FLY.io.
\emph{K1ndl1ng: a parser and normaliser for kernel rholang.}
\code{publications/campf1r3/k1ndl1ng-spec}, 2026.

\end{thebibliography}

\end{document}
